\subsection{Mencari Fungsi Asal pada Fungsi Komposisi}
\begin{definition}[Kasus 1: Mencari Fungsi Luar]
  Diketahui fungsi dalam $g(x)$ dan hasil komposisi $(f \circ g)(x)$, lalu dicari fungsi luar $f(x)$. Misalkan $g(x) = ax + b$ dan $(f \circ g)(x)$ adalah rumus yang diketahui.
\end{definition}

\begin{example}[Kasus 1]
  \hfill \\
  Diketahui:
  \[ (f \circ g)(x) = 4x^2 + 2x - 5 \]
  \[ g(x) = 2x + 1 \]
  Tentukan $f(x)$!
  
  \textbf{Penyelesaian:}
  \[
  \begin{aligned}
    f(g(x)) &= 4x^2 + 2x - 5 \\
    f(2x + 1) &= 4x^2 + 2x - 5
  \end{aligned}
  \]
  Misalkan $2x + 1 = y$, maka:
  \[
  \begin{aligned}
    2x &= y - 1 \\
    x &= \frac{y - 1}{2}
  \end{aligned}
  \]
  Substitusikan nilai $x$ tersebut ke dalam persamaan fungsi:
  \[
  \begin{aligned}
    f(y) &= 4 \left(\frac{y - 1}{2}\right)^2 + 2\left(\frac{y - 1}{2}\right) - 5 \\
    &= 4 \left(\frac{y^2 - 2y + 1}{4}\right) + (y - 1) - 5 \\
    &= (y^2 - 2y + 1) + y - 6 \\
    &= y^2 - y - 5
  \end{aligned}
  \]
  Ubah kembali variabel $y$ menjadi $x$:
  \[ f(x) = x^2 - x - 5 \]
\end{example}

\begin{definition}[Kasus 2: Mencari Fungsi Dalam]
  Diketahui fungsi luar $f(x)$ dan hasil komposisi $(f \circ g)(x)$, lalu dicari fungsi dalam $g(x)$. Masukkan $g(x)$ ke dalam rumus $f(x)$, lalu samakan dengan $(f \circ g)(x)$ yang diketahui.
\end{definition}

\begin{example}[Kasus 2]
  \hfill \\
  Diketahui:
  \[ (f \circ g)(x) = x^2 - 2x + 3 \]
  \[ f(x) = x - 4 \]
  Tentukan $g(x)$!
  
  \textbf{Penyelesaian:}
  \[
  \begin{aligned}
    f(g(x)) &= x^2 - 2x + 3 \\
    g(x) - 4 &= x^2 - 2x + 3 \\
    g(x) &= x^2 - 2x + 3 + 4 \\
    g(x) &= x^2 - 2x + 7
  \end{aligned}
  \]
\end{example}
